\par\medskip
\noindent\underline{Proposición 3.4.} \underline{Sea $P$ un esquema
en grupos liso de presentación finita sobre $S$, con fibras conexas,
y sea $Q$ un toro sobre $S$. Entonces la categoría de las
biextensiones de $(P,Q)$ por $\underline G_m$ es equivalente a la
categoría puntual; es decir,}
\begin{equation}\tag{3.4.1}
 \Biext^0(P,Q;\underline G_m)=
 \Biext^1(P,Q;\underline G_m)=0.
\end{equation}

\noindent\underline{Demostración.} Pongamos $M=D(Q)$. Este es un grupo
constante retorcido sobre $S$ cuyas fibras son isomorfas a grupos
$\mathbb Z^r$ (3.3.3). Por 3.3.2, la categoría considerada es
equivalente a la categoría de las extensiones de $P$ por $M$. Esta
categoría es \underline{rígida}: el grupo de automorfismos de todo
objeto es el grupo trivial.

%%%%%%%%%%  scan idx 271  |  folio 260  %%%%%%%%%%

Esta es la primera igualdad de (3.4.1), que también puede escribirse
\[
 \Hom(P,M)=0.
\]
Se deduce inmediatamente de que $M$ es étale sobre $S$, mientras que
$P$ tiene fibras conexas. Queda por demostrar
\[
 \Ext^1(P,M)=0,
\]
es decir, que toda extensión $E$ de $P$ por $M$ es trivial.

Por la unicidad de una escisión cuando existe, basta demostrar que $E$
es localmente trivial. Podemos, pues, reducirnos al caso
$M=\mathbb Z_S^r$, y después a $M=\mathbb Z_S$. Cuando $S$ es el
espectro de un cuerpo, la conclusión se deduce de 5.5 (i) más abajo.
(Obsérvese que no utilizaremos 3.4 antes de IX 6, por lo que no hay
ningún círculo vicioso.) Reduzcamos ahora el caso general a este caso
particular.

Basta encontrar una sección $f$ de $\sisi{P}{E}$ sobre $P$ tal que
$f(0)=0$, pues tal sección es necesariamente multiplicativa. En efecto,
\[
 f(x+y)-f(x)-f(y)=0:
\]
el miembro izquierdo define un morfismo
$P\times_S P\longrightarrow\mathbb Z_S$ que se anula sobre la sección
cero de $P\times_S P$. Por tanto, se anula idénticamente, ya que
$P\times_S P$ tiene fibras conexas (las fibras geométricas de $P$ son
conexas). La conclusión buscada se deduce ahora del lema siguiente.

\medskip
\noindent\underline{Lema 3.4.2.} \underline{Sean $S$ un esquema, $P$
un $S$-esquema liso de presentación finita con fibras geométricamente
conexas, $Z$ un grupo y $E$ un $Z_P$-torsor (para la topología fppf,
étale o fpqc; estas nociones son aquí equivalentes, cf. SGA 3 X 5.4).
Sea $e$ una sección de $P$ sobre $S$, y sea $e'$ una sección de $E$
sobre $S$ situada sobre $e$ (es decir, una sección de $e^{-1}(E)$
sobre $S$).}

%%%%%%%%%%  scan idx 272  |  folio 261  %%%%%%%%%%

\underline{El torsor $E$ es trivial si y solo si, para todo $s\in S$,
el torsor inducido $E_s$ sobre $P_s$ es trivial. En este caso existe
una única sección $f$ de $E$ sobre $P$ tal que $f\circ e=e'$.}

\noindent\underline{Demostración.} La unicidad es clara. Dos secciones
$f$ y $g$ de este tipo ``difieren'' en un morfismo
\[
 h:P\longrightarrow Z_S,
 \qquad h\circ e=1,
\]
y tal morfismo es constantemente igual a $1$, pues $P$ tiene fibras
conexas. Además, si existe una sección $f$ de $E$ sobre $P$, siempre
puede corregirse mediante un morfismo $P\longrightarrow Z_S$
procedente de un morfismo $z:S\longrightarrow Z_S$, para obtener una
sección $f'$ que satisfaga $f'\circ e=e'$. Basta elegir $z$ de modo que
\[
 f(e(t))=e'(t)z(t).
\]

Aplicando esta observación fibra a fibra, para todo $s\in S$ obtenemos
una sección bien determinada $f_s$ de $E_s$ sobre $P_s$ que lleva
$e_s$ en $e'_s$. Queda por demostrar que los $f_s$ proceden de una
sección $f$ de $E$ sobre $P$. Equivalentemente (EGA IV 17.9.3), hay
que demostrar que
\[
 U=\bigcup_{s\in S}f_s(P_s)\subset E
\]
es abierto en $E$.

Supongamos primero que $S$ es noetheriano. Entonces $U$ es
constructible. Para verlo, sea $T$, para cada $s\in S$, la clausura de
$s$, dotada de su estructura inducida reducida. Existe un entorno
abierto $T'$ de $s$ en $T$ y una sección de $E_{T'}$ sobre $P_{T'}$
que prolonga $f_s$ y es compatible con las secciones marcadas $e$ y
$e'$ (como se deduce directamente de EGA IV 8.8.2). Por el criterio de
EGA IV 1.10.1, queda por demostrar que $U$ es estable por
generización. El argumento usual que utiliza EGA II 7.1.9 reduce esto
al caso en que $S$ es el espectro de un anillo de valoración discreta.

%%%%%%%%%%  scan idx 273  |  folio 262  %%%%%%%%%%

Como $P$ es liso sobre $S$, es regular y por tanto normal; y como tiene
fibras conexas, es conexo. La sección dada de $E$ sobre la fibra
genérica de $P$ sobre $S$ se prolonga entonces de manera única a una
sección $f$ de $E$ sobre $P$ (EGA IV 18.10.19 y 18.10.20),
necesariamente compatible con $e$ y $e'$. Así, en este caso,
$U=f(P)$ es abierto. Esto prueba 3.4 cuando $S$ es noetheriano. El caso
general se reduce fácilmente a este mediante los argumentos usuales
de paso al límite de EGA IV, 8 y 9; dejamos los detalles al lector.

\medskip
\noindent\underline{Observación 3.4.3.} El lema 3.4.2 puede quedar
``aplastado'' bajo referencias eruditas, a saber, SGA 4 XV 1.15 (el
caso $n=1$) y 4.1.
